\documentclass[11pt]{article}

\usepackage[a4paper,margin=1in]{geometry}
\usepackage{amsmath,amssymb,bm}
\usepackage{booktabs}

\title{Mathematical Formulation for PINN-Based Cardiovascular Hemodynamics and Stenosis Detection}
\author{}
\date{}

\begin{document}
\maketitle

\section{Project Roadmap}

The project is organized in the following order:
\begin{enumerate}
    \item a baseline steady-flow PINN solver for a healthy straight channel,
    \item synthetic stenosed-vessel and stenosed-tree forward models,
    \item inverse recovery of stenosis location and severity, with the inverse PINN as the main detection method,
    \item richer two-dimensional stenosed-channel PINNs with known and unknown stenosis geometry.
\end{enumerate}

The reduced-order models are retained because they are fast, interpretable, and useful for creating synthetic data and sanity checks. The main stenosis-detection endpoint is not reduced-order linear algebra. It is a physics-informed inverse problem where the neural flow field and stenosis parameters are optimized together.

\section{Common Physical Assumptions and Notation}

\subsection{Fluid Model}

Blood is modeled as an incompressible Newtonian fluid in the current implementation. The density is $\rho$, the dynamic viscosity is $\mu$, and the kinematic viscosity is
\begin{equation}
    \nu = \frac{\mu}{\rho}.
\end{equation}

For the two-dimensional PINN stages, the unknown velocity-pressure field is
\begin{equation}
    \bm{q}(x,y)
    =
    \begin{bmatrix}
        u(x,y)\\
        v(x,y)\\
        p(x,y)
    \end{bmatrix},
\end{equation}
where $u$ is axial velocity, $v$ is transverse velocity, and $p$ is pressure.

The neural approximation is
\begin{equation}
    \mathcal{N}_{\theta}(x,y)
    =
    \begin{bmatrix}
        u_{\theta}(x,y)\\
        v_{\theta}(x,y)\\
        p_{\theta}(x,y)
    \end{bmatrix}.
\end{equation}

\subsection{Steady Navier--Stokes Residuals}

The steady incompressible Navier--Stokes equations are
\begin{align}
    \nabla\cdot \bm{u} &= 0,\\
    (\bm{u}\cdot\nabla)\bm{u}
    + \frac{1}{\rho}\nabla p
    - \nu \nabla^2 \bm{u}
    &= \bm{0}.
\end{align}

In two dimensions, the residuals evaluated on the neural field are
\begin{align}
    r_c
    &=
    \frac{\partial u_{\theta}}{\partial x}
    +
    \frac{\partial v_{\theta}}{\partial y},
    \label{eq:common_continuity}\\
    r_x
    &=
    u_{\theta}\frac{\partial u_{\theta}}{\partial x}
    +
    v_{\theta}\frac{\partial u_{\theta}}{\partial y}
    +
    \frac{1}{\rho}\frac{\partial p_{\theta}}{\partial x}
    -
    \nu
    \left(
        \frac{\partial^2 u_{\theta}}{\partial x^2}
        +
        \frac{\partial^2 u_{\theta}}{\partial y^2}
    \right),
    \label{eq:common_momentum_x}\\
    r_y
    &=
    u_{\theta}\frac{\partial v_{\theta}}{\partial x}
    +
    v_{\theta}\frac{\partial v_{\theta}}{\partial y}
    +
    \frac{1}{\rho}\frac{\partial p_{\theta}}{\partial y}
    -
    \nu
    \left(
        \frac{\partial^2 v_{\theta}}{\partial x^2}
        +
        \frac{\partial^2 v_{\theta}}{\partial y^2}
    \right).
    \label{eq:common_momentum_y}
\end{align}

All derivatives in the PINN losses are computed by automatic differentiation.

\subsection{Initial Conditions}

All implemented flow solvers are steady, so there is no physical time initial condition. Neural-network training starts from optimization initial values:
\begin{equation}
    \theta = \theta_0.
\end{equation}
For inverse PINNs, the stenosis parameters are also initialized:
\begin{equation}
    s=s_0,
    \qquad
    c=c_0.
\end{equation}
These are optimization initial guesses, not time-dependent physical initial conditions.

If the project later moves to pulsatile flow, the unsteady formulation would add time $t$ and require
\begin{equation}
    \bm{u}(\bm{x},0)=\bm{u}_0(\bm{x}),
    \qquad
    p(\bm{x},0)=p_0(\bm{x}),
\end{equation}
or another clinically justified initialization.

\section{Stage 1: Baseline Healthy-Channel PINN}

\subsection{Domain}

The baseline channel is
\begin{equation}
    \Omega_0
    =
    \{(x,y)\mid 0\le x\le L,\; -h\le y\le h\},
    \qquad
    h=\frac{H}{2}.
\end{equation}

\subsection{Boundary Conditions}

The inlet boundary is
\begin{equation}
    \Gamma_{\mathrm{in}}
    =
    \{(0,y)\mid -h\le y\le h\}.
\end{equation}
The prescribed Poiseuille inlet velocity is
\begin{align}
    u(0,y)
    &=
    U_{\max}
    \left(
        1-\frac{y^2}{h^2}
    \right),\\
    v(0,y) &= 0.
\end{align}

The wall boundary is
\begin{equation}
    \Gamma_{\mathrm{wall}}
    =
    \{(x,\pm h)\mid 0\le x\le L\}.
\end{equation}
The wall condition is no slip and no penetration:
\begin{align}
    u(x,\pm h)&=0,\\
    v(x,\pm h)&=0.
\end{align}

The outlet boundary is
\begin{equation}
    \Gamma_{\mathrm{out}}
    =
    \{(L,y)\mid -h\le y\le h\}.
\end{equation}
The baseline anchors transverse velocity and outlet pressure:
\begin{align}
    v(L,y)&=0,\\
    p(L,y)&=P_{\mathrm{out}}.
\end{align}

\subsection{Analytical Reference}

For validation, the exact Poiseuille solution is
\begin{align}
    u_{\mathrm{exact}}(y)
    &=
    U_{\max}
    \left(
        1-\frac{y^2}{h^2}
    \right),\\
    v_{\mathrm{exact}}(x,y)
    &=0,\\
    p_{\mathrm{exact}}(x)
    &=
    P_{\mathrm{in}}
    +
    \frac{dP}{dx}x,
\end{align}
where
\begin{equation}
    \frac{dP}{dx}
    =
    -\frac{2\rho\nu U_{\max}}{h^2}.
\end{equation}

\subsection{PINN Loss}

Let $\mathcal{C}$ be interior collocation points, $\mathcal{W}$ wall points, $\mathcal{I}$ inlet points, and $\mathcal{O}$ outlet points. The PDE residual loss is
\begin{equation}
    \mathcal{L}_{\mathrm{PDE}}
    =
    \frac{1}{|\mathcal{C}|}
    \sum_{(x_i,y_i)\in\mathcal{C}}
    \left[
        r_c(x_i,y_i)^2
        +
        r_x(x_i,y_i)^2
        +
        r_y(x_i,y_i)^2
    \right].
\end{equation}

The boundary losses are
\begin{align}
    \mathcal{L}_{\mathrm{wall}}
    &=
    \frac{1}{|\mathcal{W}|}
    \sum_{(x_i,y_i)\in\mathcal{W}}
    \left[
        u_{\theta}(x_i,y_i)^2
        +
        v_{\theta}(x_i,y_i)^2
    \right],\\
    \mathcal{L}_{\mathrm{in}}
    &=
    \frac{1}{|\mathcal{I}|}
    \sum_{(0,y_i)\in\mathcal{I}}
    \left[
        \left(
            u_{\theta}(0,y_i)
            -
            U_{\max}
            \left(
                1-\frac{y_i^2}{h^2}
            \right)
        \right)^2
        +
        v_{\theta}(0,y_i)^2
    \right],\\
    \mathcal{L}_{\mathrm{out}}
    &=
    \frac{1}{|\mathcal{O}|}
    \sum_{(L,y_i)\in\mathcal{O}}
    \left[
        v_{\theta}(L,y_i)^2
        +
        \left(
            p_{\theta}(L,y_i)-P_{\mathrm{out}}
        \right)^2
    \right].
\end{align}

The baseline optimization problem is
\begin{equation}
    \theta^{\ast}
    =
    \arg\min_{\theta}
    \mathcal{L}_{\mathrm{base}}(\theta),
\end{equation}
with
\begin{equation}
    \mathcal{L}_{\mathrm{base}}
    =
    \mathcal{L}_{\mathrm{PDE}}
    +
    10\mathcal{L}_{\mathrm{wall}}
    +
    10\mathcal{L}_{\mathrm{in}}
    +
    10\mathcal{L}_{\mathrm{out}}.
\end{equation}

\section{Stage 2: Synthetic Stenosed Forward Models}

Synthetic forward models provide controlled ground truth for stenosis severity, location, velocity, pressure drop, Reynolds number, and wall shear stress.

\subsection{Single-Vessel Geometry}

The one-dimensional vessel coordinate is
\begin{equation}
    x\in[0,L].
\end{equation}

The healthy radius is $R_0$. A Gaussian stenosis is
\begin{equation}
    R(x;s,c,w)
    =
    R_0
    \left[
        1
        -
        s
        \exp
        \left(
            -\frac{1}{2}
            \left(
                \frac{x-c}{w}
            \right)^2
        \right)
    \right],
    \label{eq:single_radius}
\end{equation}
where $s\in[0,1)$ is stenosis severity, $c$ is the stenosis center, and $w$ is the stenosis width.

The cross-sectional area is
\begin{equation}
    A(x;s,c)
    =
    \pi R(x;s,c,w)^2.
\end{equation}

\subsection{Single-Vessel Hemodynamics}

For prescribed volumetric flow rate $Q$, mass conservation gives
\begin{equation}
    Q = A(x)U(x),
\end{equation}
so the cross-section-averaged velocity is
\begin{equation}
    U(x;s,c)
    =
    \frac{Q}{A(x;s,c)}
    =
    \frac{Q}{\pi R(x;s,c,w)^2}.
\end{equation}

The local Hagen--Poiseuille pressure-gradient magnitude is
\begin{equation}
    g(x;s,c)
    =
    \frac{8\mu Q}{\pi R(x;s,c,w)^4}.
\end{equation}

The accumulated pressure drop is
\begin{equation}
    \Delta P(x;s,c)
    =
    \int_{0}^{x}
    g(\xi;s,c)\,d\xi.
\end{equation}

The pressure is propagated from the inlet:
\begin{equation}
    P(x;s,c)
    =
    P_{\mathrm{in}}
    -
    \Delta P(x;s,c).
\end{equation}

Derived quantities are
\begin{align}
    D(x;s,c) &= 2R(x;s,c,w),\\
    Re(x;s,c) &=
    \frac{\rho U(x;s,c)D(x;s,c)}{\mu},\\
    \tau_w(x;s,c) &=
    \frac{4\mu Q}{\pi R(x;s,c,w)^3},\\
    \mathcal{R}(s,c) &=
    \frac{\Delta P(L;s,c)}{Q}.
\end{align}

\subsection{Single-Vessel Boundary Conditions}

The reduced-order forward problem is steady and uses
\begin{align}
    P(0) &= P_{\mathrm{in}},\\
    Q(x) &= Q,\qquad 0\le x\le L.
\end{align}
The outlet pressure is computed:
\begin{equation}
    P(L)
    =
    P_{\mathrm{in}}
    -
    \Delta P(L).
\end{equation}

\subsection{Y-Bifurcation Tree Topology}

The tree phase uses a fixed Y topology:
\begin{equation}
    \mathcal{B}
    =
    \{b_0,b_L,b_R\},
\end{equation}
where $b_0$ is the inlet branch, $b_L$ is the left outlet branch, and $b_R$ is the right outlet branch.

Each branch has local coordinate
\begin{equation}
    x_b\in[0,L_b],
\end{equation}
healthy radius $R_{0,b}$, branch flow rate $Q_b$, and optional stenosis severity $s_b$.

\subsection{Flow Split}

The inlet flow is prescribed:
\begin{equation}
    Q_0 = Q_{\mathrm{in}}.
\end{equation}

The outlet flows are assigned by Murray-style radius weighting:
\begin{align}
    Q_L
    &=
    Q_{\mathrm{in}}
    \frac{R_{0,L}^{3}}
    {R_{0,L}^{3}+R_{0,R}^{3}},
    \\
    Q_R
    &=
    Q_{\mathrm{in}}
    \frac{R_{0,R}^{3}}
    {R_{0,L}^{3}+R_{0,R}^{3}}.
\end{align}
Thus
\begin{equation}
    Q_0 = Q_L + Q_R.
\end{equation}

\subsection{Branch Geometry and Hemodynamics}

The inlet branch is healthy:
\begin{equation}
    R_{b_0}(x_{b_0})=R_{0,b_0}.
\end{equation}

For an outlet branch $b\in\{b_L,b_R\}$,
\begin{equation}
    R_b(x_b;s_b)
    =
    R_{0,b}
    \left[
        1
        -
        s_b
        \exp
        \left(
            -\frac{1}{2}
            \left(
                \frac{x_b-c_b}{w}
            \right)^2
        \right)
    \right].
\end{equation}
Healthy branches use $s_b=0$.

For each branch,
\begin{align}
    A_b(x_b) &= \pi R_b(x_b)^2,\\
    U_b(x_b) &= \frac{Q_b}{A_b(x_b)},\\
    g_b(x_b) &= \frac{8\mu Q_b}{\pi R_b(x_b)^4},\\
    \Delta P_b(x_b) &= \int_{0}^{x_b}g_b(\xi)\,d\xi,\\
    P_b(x_b) &= P_b(0)-\Delta P_b(x_b),\\
    Re_b(x_b) &= \frac{2\rho U_b(x_b)R_b(x_b)}{\mu},\\
    \tau_{w,b}(x_b) &= \frac{4\mu Q_b}{\pi R_b(x_b)^3},\\
    \mathcal{R}_b &= \frac{\Delta P_b(L_b)}{Q_b}.
\end{align}

\subsection{Tree Boundary and Interface Conditions}

At the inlet,
\begin{align}
    P_{b_0}(0) &= P_{\mathrm{in}},\\
    Q_{b_0} &= Q_{\mathrm{in}}.
\end{align}

At the bifurcation, mass conservation and pressure continuity are
\begin{align}
    Q_{b_0} &= Q_{b_L}+Q_{b_R},\\
    P_{b_L}(0)
    =
    P_{b_R}(0)
    =
    P_{b_0}(L_{b_0})
    &=
    P_{\mathrm{bif}}.
\end{align}

Outlet pressures are computed by branch pressure drops:
\begin{align}
    P_{b_L}(L_{b_L})
    &=
    P_{\mathrm{bif}}
    -
    \Delta P_{b_L}(L_{b_L}),
    \\
    P_{b_R}(L_{b_R})
    &=
    P_{\mathrm{bif}}
    -
    \Delta P_{b_R}(L_{b_R}).
\end{align}

\section{Stage 3: Inverse Recovery of Stenosis}

\subsection{Reduced-Order Single-Vessel Inverse Baseline}

The reduced-order inverse baseline estimates
\begin{equation}
    \lambda=(s,c).
\end{equation}

Unconstrained raw variables are mapped into physical bounds:
\begin{align}
    s(\alpha_s)
    &=
    s_{\min}
    +
    (s_{\max}-s_{\min})\sigma(\alpha_s),\\
    c(\alpha_c)
    &=
    c_{\min}
    +
    (c_{\max}-c_{\min})\sigma(\alpha_c),
\end{align}
where
\begin{equation}
    \sigma(z)=\frac{1}{1+e^{-z}}.
\end{equation}

Sparse measurements are sampled at
\begin{equation}
    \mathcal{S}=\{x_j\}_{j=1}^{N_s}.
\end{equation}
The observations are
\begin{align}
    U_j^{\mathrm{obs}}
    &=
    U(x_j;s_{\mathrm{true}},c_{\mathrm{true}})
    +
    \epsilon_j^U,\\
    \Delta P_j^{\mathrm{obs}}
    &=
    \Delta P(x_j;s_{\mathrm{true}},c_{\mathrm{true}})
    +
    \epsilon_j^P.
\end{align}

For candidate $(s,c)$,
\begin{align}
    U_j(s,c)
    &=
    \frac{Q}{\pi R(x_j;s,c,w)^2},\\
    \Delta P_j(s,c)
    &=
    \int_{0}^{x_j}
    \frac{8\mu Q}{\pi R(\xi;s,c,w)^4}\,d\xi.
\end{align}

The inverse loss is
\begin{equation}
    \mathcal{L}_{\mathrm{RO,1D}}(s,c)
    =
    \frac{1}{N_s}
    \sum_{j=1}^{N_s}
    \left[
        \left(
            \frac{U_j(s,c)-U_j^{\mathrm{obs}}}
            {U_{\mathrm{scale}}}
        \right)^2
        +
        \left(
            \frac{\Delta P_j(s,c)-\Delta P_j^{\mathrm{obs}}}
            {P_{\mathrm{scale}}}
        \right)^2
    \right].
\end{equation}

The reduced-order estimate is
\begin{equation}
    (s^{\ast},c^{\ast})
    =
    \arg\min_{s,c}
    \mathcal{L}_{\mathrm{RO,1D}}(s,c).
\end{equation}

\subsection{Reduced-Order Tree Inverse Baseline}

For the Y tree, the inverse baseline estimates outlet severities
\begin{equation}
    \lambda_{\mathrm{tree}}=(s_L,s_R).
\end{equation}

The branch centers are fixed in the current implementation:
\begin{equation}
    c_L=\frac{1}{2}L_{b_L},
    \qquad
    c_R=\frac{1}{2}L_{b_R}.
\end{equation}

For each outlet branch $b\in\{L,R\}$, sparse observations are
\begin{equation}
    \mathcal{S}_b=\{x_{b,j}\}_{j=1}^{N_b},
\end{equation}
with
\begin{align}
    U_{b,j}^{\mathrm{obs}}
    &=
    U_b(x_{b,j};s_b^{\mathrm{true}})
    +
    \epsilon_{b,j}^{U},\\
    \Delta P_{b,j}^{\mathrm{obs}}
    &=
    \Delta P_b(x_{b,j};s_b^{\mathrm{true}})
    +
    \epsilon_{b,j}^{P}.
\end{align}

The branch loss is
\begin{equation}
    \mathcal{L}_b(s_b)
    =
    \frac{1}{N_b}
    \sum_{j=1}^{N_b}
    \left[
        \left(
            \frac{U_b(x_{b,j};s_b)-U_{b,j}^{\mathrm{obs}}}
            {U_{b,\mathrm{scale}}}
        \right)^2
        +
        \left(
            \frac{\Delta P_b(x_{b,j};s_b)-\Delta P_{b,j}^{\mathrm{obs}}}
            {P_{b,\mathrm{scale}}}
        \right)^2
    \right].
\end{equation}

The tree inverse loss is
\begin{equation}
    \mathcal{L}_{\mathrm{RO,tree}}(s_L,s_R)
    =
    \mathcal{L}_L(s_L)+\mathcal{L}_R(s_R).
\end{equation}

The recovered outlet severities are
\begin{equation}
    (s_L^{\ast},s_R^{\ast})
    =
    \arg\min_{s_L,s_R}
    \mathcal{L}_{\mathrm{RO,tree}}(s_L,s_R).
\end{equation}

The classified stenosed branch is
\begin{equation}
    b^{\ast}
    =
    \begin{cases}
        b_L, & s_L^{\ast}\ge s_R^{\ast},\\
        b_R, & s_R^{\ast}>s_L^{\ast},
    \end{cases}
\end{equation}
and the recovered severity is
\begin{equation}
    s^{\ast}=\max(s_L^{\ast},s_R^{\ast}).
\end{equation}

\subsection{Main Method: Inverse PINN Stenosis Detection}

The inverse PINN is the main stenosis-detection formulation. It estimates the flow field and the unknown blockage parameters in one physics-constrained optimization problem.

\subsubsection{Unknown Stenosed Domain}

The stenosed channel half-height is
\begin{equation}
    h(x;s,c)
    =
    h_0
    \left[
        1
        -
        s
        \exp
        \left(
            -\frac{1}{2}
            \left(
                \frac{x-c}{w}
            \right)^2
        \right)
    \right],
    \qquad
    h_0=\frac{H}{2}.
    \label{eq:inverse_height}
\end{equation}

The domain depends on the unknown stenosis parameters:
\begin{equation}
    \Omega(s,c)
    =
    \{(x,y)\mid 0\le x\le L,\; -h(x;s,c)\le y\le h(x;s,c)\}.
\end{equation}

The trainable unknowns are
\begin{equation}
    \Lambda=(\theta,s,c).
\end{equation}

\subsubsection{PINN Residual}

At each training step, collocation points are sampled inside $\Omega(s,c)$. The PDE residuals are the same Navier--Stokes residuals from Equations~\eqref{eq:common_continuity}--\eqref{eq:common_momentum_y}.

The inverse-PINN PDE loss is
\begin{equation}
    \mathcal{L}_{\mathrm{PDE}}^{\mathrm{inv}}
    =
    \frac{1}{|\mathcal{C}(s,c)|}
    \sum_{(x_i,y_i)\in\mathcal{C}(s,c)}
    \left[
        4r_c(x_i,y_i)^2
        +
        r_x(x_i,y_i)^2
        +
        r_y(x_i,y_i)^2
    \right].
\end{equation}

\subsubsection{Inverse-PINN Boundary Conditions}

The inferred wall is
\begin{equation}
    \Gamma_{\mathrm{wall}}(s,c)
    =
    \{(x,\pm h(x;s,c))\mid 0\le x\le L\}.
\end{equation}
No slip and no penetration give
\begin{align}
    u_{\theta}(x,\pm h(x;s,c))&=0,\\
    v_{\theta}(x,\pm h(x;s,c))&=0.
\end{align}

At the inlet,
\begin{align}
    u_{\theta}(0,y)
    &=
    U_{\max}
    \left(
        1-\frac{y^2}{h_0^2}
    \right),\\
    v_{\theta}(0,y)&=0,\\
    p_{\theta}(0,y)&=P_{\mathrm{in}}.
\end{align}

At the outlet,
\begin{align}
    v_{\theta}(L,y)&=0,\\
    p_{\theta}(L,y)&=P_{\mathrm{out}}.
\end{align}

The boundary loss is
\begin{equation}
    \mathcal{L}_{\mathrm{BC}}^{\mathrm{inv}}
    =
    \mathcal{L}_{\mathrm{wall}}
    +
    \mathcal{L}_{\mathrm{in}}
    +
    \mathcal{L}_{\mathrm{out}}.
\end{equation}

\subsubsection{Sparse Flow, Pressure, and Wall Observations}

Sparse flow-pressure observations are
\begin{equation}
    \mathcal{D}_f
    =
    \{
    (x_j,y_j,u_j^{\mathrm{obs}},v_j^{\mathrm{obs}},p_j^{\mathrm{obs}})
    \}_{j=1}^{N_f}.
\end{equation}

The data loss is
\begin{equation}
    \mathcal{L}_{\mathrm{data}}^{\mathrm{inv}}
    =
    \frac{1}{N_f}
    \sum_{j=1}^{N_f}
    \left[
        \left(
            \frac{u_{\theta}(x_j,y_j)-u_j^{\mathrm{obs}}}
            {U_{\mathrm{scale}}}
        \right)^2
        +
        \left(
            \frac{v_{\theta}(x_j,y_j)-v_j^{\mathrm{obs}}}
            {U_{\mathrm{scale}}}
        \right)^2
        +
        \left(
            \frac{p_{\theta}(x_j,y_j)-p_j^{\mathrm{obs}}}
            {P_{\mathrm{scale}}}
        \right)^2
    \right].
\end{equation}

Sparse anatomical wall observations are
\begin{equation}
    \mathcal{D}_g
    =
    \{(x_k,h_k^{\mathrm{obs}})\}_{k=1}^{N_g}.
\end{equation}
These make the stenosis geometry identifiable, as would wall contours from imaging. The geometry loss is
\begin{equation}
    \mathcal{L}_{\mathrm{geom}}
    =
    \frac{1}{N_g}
    \sum_{k=1}^{N_g}
    \left(
        \frac{h(x_k;s,c)-h_k^{\mathrm{obs}}}{h_0}
    \right)^2.
\end{equation}

\subsubsection{Joint Inverse-PINN Objective}

The implemented inverse-PINN loss is
\begin{equation}
    \mathcal{L}_{\mathrm{invPINN}}(\theta,s,c)
    =
    \mathcal{L}_{\mathrm{PDE}}^{\mathrm{inv}}
    +
    10\mathcal{L}_{\mathrm{BC}}^{\mathrm{inv}}
    +
    5\mathcal{L}_{\mathrm{data}}^{\mathrm{inv}}
    +
    200\mathcal{L}_{\mathrm{geom}}.
\end{equation}

The final recovered stenosis parameters are obtained from
\begin{equation}
    (\theta^{\ast},s^{\ast},c^{\ast})
    =
    \arg\min_{\theta,s,c}
    \mathcal{L}_{\mathrm{invPINN}}(\theta,s,c).
\end{equation}

This is the central detection step: the stenosis parameters are recovered while the neural velocity-pressure field obeys the Navier--Stokes residual, boundary conditions, sparse sensor data, and sparse vessel-wall data.

\section{Stage 4: Richer 2D Stenosed-Channel PINNs}

\subsection{Forward 2D Stenosed-Channel PINN}

The forward 2D stenosed-channel PINN uses the same geometry as Equation~\eqref{eq:inverse_height}, but the stenosis parameters are known:
\begin{equation}
    h_{\mathrm{known}}(x)
    =
    h(x;s_{\mathrm{given}},c_{\mathrm{given}}).
\end{equation}
The domain is
\begin{equation}
    \Omega_{\mathrm{sten}}
    =
    \{(x,y)\mid 0\le x\le L,\; -h_{\mathrm{known}}(x)\le y\le h_{\mathrm{known}}(x)\}.
\end{equation}

The network $\mathcal{N}_{\theta}(x,y)$ solves the same steady Navier--Stokes residuals. The wall, inlet, and outlet conditions are
\begin{align}
    u_{\theta}(x,\pm h_{\mathrm{known}}(x))&=0,\\
    v_{\theta}(x,\pm h_{\mathrm{known}}(x))&=0,\\
    u_{\theta}(0,y)
    &=
    U_{\max}
    \left(
        1-\frac{y^2}{h_0^2}
    \right),\\
    v_{\theta}(0,y)&=0,\\
    p_{\theta}(0,y)&=P_{\mathrm{in}},\\
    v_{\theta}(L,y)&=0,\\
    p_{\theta}(L,y)&=P_{\mathrm{out}}.
\end{align}

To improve training in the curved stenosed region, collocation points are not sampled only uniformly. A fraction of points is concentrated near the stenosis center $c_{\mathrm{given}}$ and a fraction is concentrated near the walls:
\begin{equation}
    x_i
    \sim
    (1-\alpha_{\ell})\mathcal{U}(0,L)
    +
    \alpha_{\ell}
    \mathcal{N}(c_{\mathrm{given}},w^2),
\end{equation}
with clipping to $[0,L]$. The transverse coordinate is sampled as
\begin{equation}
    y_i=\eta_i h_{\mathrm{known}}(x_i),
    \qquad
    \eta_i\in[-1,1],
\end{equation}
with additional samples biased toward $|\eta_i|\approx 1$ to better resolve wall behavior.

The upgraded forward PINN also uses a weak synthetic-reference term. The reference is not the final solver; it is a training scaffold from the reduced-order stenosed model. For sampled points $\mathcal{D}_{\mathrm{ref}}$, define
\begin{equation}
    \begin{aligned}
        \mathcal{L}_{\mathrm{ref}}
        &=
        \frac{1}{|\mathcal{D}_{\mathrm{ref}}|}
        \sum_{(x_j,y_j)\in\mathcal{D}_{\mathrm{ref}}}
        \bigg[
            \left(
                \frac{u_{\theta}(x_j,y_j)-U_{\mathrm{ref}}(x_j,y_j)}
                {U_{\max}}
            \right)^2 \\
        &\qquad
            +
            \left(
                \frac{v_{\theta}(x_j,y_j)}
                {U_{\max}}
            \right)^2
            +
            \left(
                \frac{p_{\theta}(x_j,y_j)-P_{\mathrm{ref}}(x_j)}
                {P_{\mathrm{in}}-P_{\mathrm{out}}}
            \right)^2
        \bigg].
    \end{aligned}
\end{equation}

The forward stenosed PINN objective is
\begin{equation}
    \theta^{\ast}
    =
    \arg\min_{\theta}
    \left[
        6\mathcal{L}_{c}
        +
        4\mathcal{L}_{m}
        +
        10\mathcal{L}_{\mathrm{wall}}
        +
        10\mathcal{L}_{\mathrm{in}}
        +
        10\mathcal{L}_{\mathrm{out}}
        +
        0.5\mathcal{L}_{\mathrm{ref}}
    \right].
\end{equation}
Here $\mathcal{L}_{c}$ is the continuity residual loss and $\mathcal{L}_{m}$ is the combined $x$- and $y$-momentum residual loss. This reweighting improves the harder curved-geometry residuals while keeping the solution physics-informed.

\section{Complete Solution Order}

The full project can be summarized as the following mathematical pipeline:
\begin{enumerate}
    \item Learn the healthy-channel PINN:
    \begin{equation}
        \theta_{\mathrm{base}}^{\ast}
        =
        \arg\min_{\theta}
        \mathcal{L}_{\mathrm{base}}(\theta).
    \end{equation}
    \item Generate synthetic stenosed single-vessel and tree fields from $R(x;s,c,w)$ and branch laws:
    \begin{equation}
        \{R,A,U,P,\Delta P,Re,\tau_w,\mathcal{R}\}
        =
        \mathcal{F}_{\mathrm{RO}}(s,c,Q).
    \end{equation}
    \item Use reduced-order inverse models only as fast baselines:
    \begin{equation}
        \lambda_{\mathrm{RO}}^{\ast}
        =
        \arg\min_{\lambda}
        \mathcal{L}_{\mathrm{RO}}(\lambda).
    \end{equation}
    \item Perform the main stenosis detection with the inverse PINN:
    \begin{equation}
        (\theta^{\ast},s^{\ast},c^{\ast})
        =
        \arg\min_{\theta,s,c}
        \mathcal{L}_{\mathrm{invPINN}}(\theta,s,c).
    \end{equation}
    \item Verify the richer 2D stenosed-channel forward PINN with known geometry and the inverse 2D PINN with unknown stenosis parameters.
\end{enumerate}

\end{document}
